\section{Khảo sát nghiệp vụ và nguồn dữ liệu}\label{sec:survey}

\subsection{Quy trình nghiệp vụ đặt phòng}

Quy trình đặt phòng của Hotelia có ba nhóm người tham gia, mỗi nhóm có mục tiêu khác nhau. Khách muốn tìm nhanh một nơi ở phù hợp và chắc chắn có phòng khi đến. Chủ khách sạn muốn bán được phòng, nắm được ai sẽ đến vào ngày nào và ghi nhận khách nhận/trả phòng. Quản trị viên muốn bảo đảm chỉ những khách sạn và chủ khách sạn hợp lệ xuất hiện trên nền tảng, và có công cụ xử lý khi phát sinh vấn đề.

\Cref{fig:act-overall} mô tả quy trình tổng thể bằng biểu đồ hoạt động có làn (swimlane). Quy trình bắt đầu khi khách tìm kiếm, đi qua việc chọn loại phòng và giữ phòng, phân nhánh theo phương thức thanh toán, rồi kết thúc bằng việc chủ khách sạn ghi nhận nhận phòng, trả phòng và khách gửi đánh giá. Hai điểm đáng chú ý trong biểu đồ là: (i) giữ phòng có thời hạn, nên mọi nhánh chờ đều có lối ra ``hết hạn''; và (ii) với payOS, trạng thái đặt phòng chỉ đổi khi máy chủ nhận được webhook đã xác minh, không phải khi khách quay lại ứng dụng.

\begin{figure}[H]
\centering
\diagram{act-overall}
\caption{Biểu đồ hoạt động tổng thể của quy trình đặt phòng}\label{fig:act-overall}
\end{figure}

\subsection{Nguồn dữ liệu của hệ thống}\label{sec:data-sources}

Việc phân biệt rõ dữ liệu thật và dữ liệu mô phỏng là điều kiện để đánh giá trung thực, đặc biệt với phần gợi ý. \Cref{tab:data-sources} liệt kê nguồn gốc của từng nhóm dữ liệu trong cơ sở dữ liệu vận hành và trong đường ống huấn luyện.

\begin{rtable}
\begin{longtable}{|L{3.4cm}|L{6.3cm}|L{5.0cm}|}
\caption{Nguồn gốc dữ liệu của Hotelia}\label{tab:data-sources}\\
\hline
\thead{Nhóm dữ liệu} & \thead{Nguồn và cách tạo} & \thead{Tính chất} \\ \hline
\endfirsthead
\hline
\thead{Nhóm dữ liệu} & \thead{Nguồn và cách tạo} & \thead{Tính chất} \\ \hline
\endhead
20 khách sạn ban đầu & Tệp \code{real-hotels.v1.json}: tên, địa chỉ, mô tả, hạng sao và tiện nghi được đối chiếu với website chính thức của khách sạn ngày 25/09/2026; 80 ảnh (20 ảnh khách sạn, 60 ảnh phòng) có mã băm SHA-256 trong \code{verified-media.v1.json}. & Thông tin mô tả đã kiểm chứng một phần; giá, sức chứa, số phòng và chủ khách sạn là dữ liệu ứng dụng, không phải giá bán thật. \\ \hline
80 khách sạn mở rộng & Tệp \code{explore-hotels.v2.json}: 30 tên có nguồn từ nhà vận hành; 50 tên và tọa độ lấy từ bộ dữ liệu TripAdvisor Vietnam Hotel Reviews (CC BY 4.0) \cite{tripadvisorvn}. Kịch bản \code{activate-operational-catalog.cjs} sinh 80 chủ khách sạn với email \code{@owners.hotelia.invalid} (không đăng nhập được) và 320 loại phòng có giá, tồn kho tất định. & Danh mục vận hành mô phỏng. Không chứng minh khách sạn hiện đang hoạt động hay hợp tác với Hotelia. \\ \hline
Điểm đánh giá biên tập & Trường \code{demoRating}, \code{demoReviewCount}, \code{demoReviews}; ứng dụng hiển thị là ``Điểm Hotelia'' kèm ghi chú ``không phải đánh giá của khách đã lưu trú''. & Dữ liệu minh họa, tách khỏi đánh giá đã xác minh (\code{averageRating}, \code{reviewCount}). \\ \hline
Đánh giá đã xác minh & Do khách tạo sau khi đặt phòng chuyển sang trạng thái đã trả phòng. & Dữ liệu thật của hệ thống; trên môi trường vận hành ngày 24/10/2026 số lượng bằng 0. \\ \hline
Thời tiết & Open-Meteo, gọi trực tiếp từ ứng dụng iOS \cite{openmeteo}. & Dữ liệu thật, không lưu trên máy chủ. \\ \hline
Tương tác gợi ý & Bộ sưu tập \code{recommendation\_impressions} và \code{recommendation\_views}, ghi khi ứng dụng hiển thị hàng gợi ý và khi khách mở chi tiết. & Dữ liệu thật nhưng chưa đủ lớn để huấn luyện; tự xóa sau 90 ngày. \\ \hline
Dữ liệu huấn luyện LightGBM & Thuộc tính khách sạn lấy từ ảnh chụp danh mục của API (\code{hotelia-catalog-snapshot.json}); lịch sử người dùng và lựa chọn được \emph{sinh} bằng mô phỏng (\cref{sec:rec-design}). & Dữ liệu trước phát hành; chỉ dùng để kiểm chứng đường ống, không ghi vào MongoDB. \\ \hline
Trivago RecSys 2019 & Bộ dữ liệu công khai của cuộc thi \cite{adamczak2019}. & Chỉ dùng cho thí nghiệm phụ; không ánh xạ được sang khách sạn Hotelia. \\ \hline
\end{longtable}
\end{rtable}

\section{Phân tích yêu cầu}\label{sec:requirements}

\subsection{Yêu cầu chức năng}

\Cref{tab:fr} liệt kê các yêu cầu chức năng, mức ưu tiên (C: cao, T: trung bình, Th: thấp) và trạng thái triển khai được đối chiếu với mã nguồn. Ba trạng thái được dùng là ``Đã có'' (có mã nguồn và đã kiểm thử hoặc kiểm tra thủ công), ``Một phần'' (có mã nguồn nhưng phụ thuộc cấu hình bên ngoài hoặc thiếu một nhánh) và ``Chưa'' (chưa triển khai).

\begin{rtable}
\begin{longtable}{|C{1.25cm}|L{7.3cm}|C{1.0cm}|C{1.5cm}|L{3.6cm}|}
\caption{Yêu cầu chức năng và trạng thái triển khai}\label{tab:fr}\\
\hline
\thead{Mã} & \thead{Yêu cầu} & \thead{Ưu tiên} & \thead{Trạng thái} & \thead{Ghi chú} \\ \hline
\endfirsthead
\hline
\thead{Mã} & \thead{Yêu cầu} & \thead{Ưu tiên} & \thead{Trạng thái} & \thead{Ghi chú} \\ \hline
\endhead
\multicolumn{5}{|l|}{\textit{Tài khoản và xác thực}} \\ \hline
FR-01 & Đăng ký bằng email, mật khẩu, họ tên, số điện thoại & C & Đã có & \\ \hline
FR-02 & Đăng nhập email -- mật khẩu, khóa 15 phút sau 5 lần sai & C & Đã có & \\ \hline
FR-03 & Đăng nhập Google bằng ID token đã xác minh & C & Đã có & Cần mã ứng dụng Google \\ \hline
FR-04 & Đăng nhập Apple với nonce & T & Một phần & Cần Apple Developer Program \\ \hline
FR-05 & Làm mới phiên bằng refresh token xoay vòng; đăng xuất thu hồi phiên trên máy chủ & C & Đã có & \\ \hline
FR-06 & Quên mật khẩu, đặt lại bằng mã có hạn 15 phút & T & Một phần & Cần cấu hình dịch vụ gửi email \\ \hline
FR-07 & Đổi mật khẩu, vô hiệu hóa các phiên cũ & T & Đã có & \\ \hline
FR-08 & Xem và cập nhật hồ sơ & T & Đã có & \\ \hline
\multicolumn{5}{|l|}{\textit{Khám phá và gợi ý}} \\ \hline
FR-09 & Tìm kiếm theo điểm đến, ngày, số khách, khoảng giá, điểm tối thiểu, tiện nghi; sắp xếp theo giá hoặc điểm & C & Đã có & \\ \hline
FR-10 & Xem chi tiết khách sạn, ảnh, tiện nghi, chính sách giờ nhận/trả phòng, đánh giá & C & Đã có & \\ \hline
FR-11 & Hàng ``Đề xuất cho bạn'' cá nhân hóa theo lịch sử xem và đặt phòng & C & Đã có & LightGBM v3 \\ \hline
FR-12 & Gợi ý khám phá khi chưa có lịch sử, đa dạng thành phố & T & Đã có & \\ \hline
FR-13 & Xóa lịch sử xem và dữ liệu gợi ý của chính mình & T & Đã có & \\ \hline
FR-14 & Thẻ thời tiết cho các thành phố lớn & Th & Đã có & 8 thành phố \\ \hline
FR-15 & Lưu khách sạn yêu thích & Th & Đã có & \\ \hline
\multicolumn{5}{|l|}{\textit{Đặt phòng và thanh toán}} \\ \hline
FR-16 & Hiển thị số phòng còn trống của từng loại phòng cho khoảng ngày đã chọn & C & Đã có & Theo từng đêm \\ \hline
FR-17 & Tạo đặt phòng, giữ phòng 15 phút, không bán trùng & C & Đã có & \\ \hline
FR-18 & Gửi lại yêu cầu đặt phòng không tạo bản sao (Idempotency-Key) & C & Đã có & \\ \hline
FR-19 & Xác nhận trả tại khách sạn & C & Đã có & \\ \hline
FR-20 & Thanh toán qua payOS, xác nhận bằng webhook ký số & C & Đã có & Cần khóa payOS \\ \hline
FR-21 & Hủy đặt phòng chưa thanh toán; gửi yêu cầu hủy với đặt phòng đã xác nhận & C & Đã có & \\ \hline
FR-22 & Hoàn tiền tự động qua cổng thanh toán & T & Chưa & Đối soát thủ công \\ \hline
FR-23 & Lịch sử đặt phòng theo nhóm (sắp tới, đã xong, đã hủy) & C & Đã có & \\ \hline
FR-24 & Đánh giá một lần cho mỗi khách sạn sau khi trả phòng & T & Đã có & \\ \hline
FR-25 & Thông báo trong ứng dụng khi giữ phòng, thanh toán & Th & Đã có & Chưa có thông báo đẩy \\ \hline
\multicolumn{5}{|l|}{\textit{Chủ khách sạn}} \\ \hline
FR-26 & Đăng ký làm chủ khách sạn, chờ quản trị viên duyệt & C & Đã có & \\ \hline
FR-27 & Tạo, sửa, gửi duyệt, ngừng khách sạn & C & Đã có & \\ \hline
FR-28 & Quản lý loại phòng; không cho giảm số phòng dưới số đã bán trong tương lai & C & Đã có & \\ \hline
FR-29 & Xem đặt phòng; nhận phòng, trả phòng; xác nhận hủy & C & Đã có & \\ \hline
FR-30 & Bảng điều khiển doanh thu theo tháng, tỷ lệ phòng đang có khách & T & Đã có & \\ \hline
FR-31 & Phản hồi đánh giá của khách & T & Đã có & \\ \hline
FR-32 & Chuyển đặt phòng đã trả phòng sang ``hoàn thành'' & Th & Chưa & Trạng thái đã khai báo \\ \hline
\multicolumn{5}{|l|}{\textit{Quản trị viên}} \\ \hline
FR-33 & Duyệt hoặc từ chối chủ khách sạn và khách sạn & C & Đã có & \\ \hline
FR-34 & Tìm kiếm, khóa, mở khóa người dùng & C & Đã có & \\ \hline
FR-35 & Ẩn đánh giá vi phạm, tính lại điểm trung bình & T & Đã có & \\ \hline
FR-36 & Xem thống kê, đặt phòng, giao dịch thanh toán, nhật ký thao tác & T & Đã có & 200 bản ghi gần nhất \\ \hline
FR-37 & Tiếp nhận và trả lời yêu cầu hỗ trợ & Th & Đã có & \\ \hline
\end{longtable}
\end{rtable}

\subsection{Yêu cầu phi chức năng}

Yêu cầu phi chức năng được phân nhóm theo mô hình chất lượng sản phẩm ISO/IEC 25010 \cite{iso25010}. \Cref{tab:nfr} nêu từng yêu cầu cùng cách hệ thống đáp ứng. Đồ án chưa thực hiện kiểm thử tải nên không đặt chỉ tiêu định lượng về thông lượng hay độ trễ. Thay vào đó, các yêu cầu hiệu năng được diễn đạt dưới dạng giới hạn thiết kế có thể kiểm chứng trong mã nguồn.

\begin{rtable}
\begin{longtable}{|C{1.4cm}|L{2.6cm}|L{5.6cm}|L{5.0cm}|}
\caption{Yêu cầu phi chức năng}\label{tab:nfr}\\
\hline
\thead{Mã} & \thead{Đặc tính} & \thead{Yêu cầu} & \thead{Cách đáp ứng} \\ \hline
\endfirsthead
\hline
\thead{Mã} & \thead{Đặc tính} & \thead{Yêu cầu} & \thead{Cách đáp ứng} \\ \hline
\endhead
NFR-01 & Tính đúng & Không bán quá số phòng của một loại phòng trong bất kỳ đêm nào & Giao dịch MongoDB, khóa theo loại phòng, đếm theo đêm (\cref{sec:booking-rules}) \\ \hline
NFR-02 & Tính đúng & Gửi lại cùng yêu cầu đặt phòng không tạo bản ghi mới & Idempotency-Key, chỉ mục duy nhất một phần \\ \hline
NFR-03 & Bảo mật & Mật khẩu không lưu dạng rõ; refresh token chỉ lưu bản băm & bcrypt (chi phí 12) cho mật khẩu, SHA-256 cho refresh token \\ \hline
NFR-04 & Bảo mật & Access token hết hạn sau 15 phút; refresh token không dùng được thay access token & Claim \code{typ}, kiểm tra trong JwtStrategy \\ \hline
NFR-05 & Bảo mật & Mọi đầu vào được kiểm tra kiểu và độ dài; trường lạ bị từ chối & ValidationPipe với \code{whitelist}, \code{forbidNonWhitelisted} \\ \hline
NFR-06 & Bảo mật & Giới hạn tần suất 100 yêu cầu/phút; CORS chỉ cho nguồn được khai báo ở môi trường vận hành & ThrottlerGuard; kiểm tra cấu hình khi khởi động \\ \hline
NFR-07 & Bảo mật & Không tin trạng thái thanh toán do máy khách gửi & Chỉ webhook có chữ ký hợp lệ và đúng số tiền \\ \hline
NFR-08 & Khả năng truy vết & Ghi lại người thực hiện các thao tác thay đổi dữ liệu & AuditInterceptor, bộ sưu tập \code{audit\_logs} \\ \hline
NFR-09 & Tin cậy & Gợi ý vẫn hoạt động khi mô hình lỗi hoặc thiếu & Bộ xếp hạng theo nội dung làm dự phòng \\ \hline
NFR-10 & Tin cậy & Gọi dịch vụ ngoài có thời gian chờ & payOS 10 giây (2 lần thử lại), Open-Meteo 12 giây, email 10 giây \\ \hline
NFR-11 & Hiệu năng & Danh sách có phân trang, tối đa 100 phần tử mỗi trang; truy vấn chính có chỉ mục & \code{limit} $\le 100$; chỉ mục trên \code{status}, \code{city}, \code{roomTypeId}, \code{checkIn}, \code{checkOut} \\ \hline
NFR-12 & Khả dụng & Có điểm kiểm tra sức khỏe cho nền tảng triển khai & \code{GET /health} \\ \hline
NFR-13 & Bản địa hóa & Ứng dụng iOS hỗ trợ tiếng Việt và tiếng Anh, đổi ngôn ngữ không cần cài lại & String Catalog, \code{AppLanguage} \\ \hline
NFR-14 & Quyền riêng tư & Dữ liệu tương tác gợi ý có thời hạn và người dùng tự xóa được & Chỉ mục TTL 90 ngày; \code{DELETE /recommendations/views} \\ \hline
NFR-15 & Bảo trì & Quy tắc nghiệp vụ quan trọng có kiểm thử tự động & 50 ca Jest, 80 ca pytest, kịch bản MongoDB thật \\ \hline
NFR-16 & Khả chuyển & API chạy trong container; cấu hình qua biến môi trường & Dockerfile hai giai đoạn, \code{render.yaml} \\ \hline
\end{longtable}
\end{rtable}

\section{Tác nhân của hệ thống}\label{sec:actors}

Hệ thống có ba tác nhân chính là con người và ba tác nhân phụ là hệ thống bên ngoài (\cref{tab:actors}). Trong đó, nhà cung cấp danh tính (Google, Apple) được bổ sung so với danh sách ban đầu, vì nó tham gia trực tiếp vào ca sử dụng đăng nhập.

\begin{table}[H]
\centering
\caption{Tác nhân của hệ thống}\label{tab:actors}
\begin{rtable}
\begin{tabular}{|L{3.2cm}|C{1.6cm}|L{9.9cm}|}
\hline
\thead{Tác nhân} & \thead{Loại} & \thead{Mô tả} \\ \hline
Khách (Guest/User) & Chính & Người dùng ứng dụng iOS. Khi chưa đăng nhập (Guest) có thể xem khách sạn và gợi ý khám phá; khi đã đăng nhập (vai trò USER) có thể đặt phòng, thanh toán, hủy, đánh giá. \\ \hline
Chủ khách sạn (Owner) & Chính & Người dùng có vai trò OWNER sau khi được duyệt. Quản lý khách sạn, loại phòng, đặt phòng, nhận/trả phòng, phản hồi đánh giá, xem doanh thu qua bảng điều khiển web hoặc màn hình quản lý trên iOS. \\ \hline
Quản trị viên (Admin) & Chính & Người dùng có vai trò ADMIN. Duyệt chủ khách sạn và khách sạn, quản lý người dùng, kiểm duyệt đánh giá, xem thanh toán và nhật ký thao tác. \\ \hline
Cổng thanh toán (payOS) & Phụ & Tạo liên kết thanh toán VietQR và gửi webhook ký số khi giao dịch hoàn tất. \\ \hline
Dịch vụ thời tiết (Open-Meteo) & Phụ & Cung cấp thời tiết hiện tại và dự báo theo tọa độ cho ứng dụng iOS. \\ \hline
Nhà cung cấp danh tính (Google, Apple) & Phụ & Xác thực người dùng và phát hành ID token cho đăng nhập liên kết. \\ \hline
\end{tabular}
\end{rtable}
\end{table}

\section{Mô hình ca sử dụng}\label{sec:usecases}

\subsection{Danh sách ca sử dụng}

\Cref{tab:uc-list} liệt kê 37 ca sử dụng, chia theo tác nhân. Mã ca sử dụng được dùng lại trong các biểu đồ và bảng đặc tả phía sau.

\begin{rtable}
\begin{longtable}{|C{1.2cm}|L{5.4cm}|L{2.5cm}|L{5.6cm}|}
\caption{Danh sách ca sử dụng}\label{tab:uc-list}\\
\hline
\thead{Mã} & \thead{Tên ca sử dụng} & \thead{Tác nhân} & \thead{Mô tả ngắn} \\ \hline
\endfirsthead
\hline
\thead{Mã} & \thead{Tên ca sử dụng} & \thead{Tác nhân} & \thead{Mô tả ngắn} \\ \hline
\endhead
UC01 & Đăng ký tài khoản & Khách & Tạo tài khoản USER bằng email và mật khẩu \\ \hline
UC02 & Đăng nhập email -- mật khẩu & Khách, Owner, Admin & Nhận cặp access/refresh token \\ \hline
UC03 & Đăng nhập Google/Apple & Khách, IdP & Đổi ID token đã xác minh lấy phiên Hotelia \\ \hline
UC04 & Làm mới phiên, đăng xuất & Mọi người dùng & Xoay vòng refresh token; thu hồi khi đăng xuất \\ \hline
UC05 & Quên và đặt lại mật khẩu & Khách & Nhận mã qua email, đặt mật khẩu mới \\ \hline
UC06 & Đổi mật khẩu & Mọi người dùng & Vô hiệu hóa token cũ \\ \hline
UC07 & Xem, cập nhật hồ sơ & Mọi người dùng & Họ tên, số điện thoại, ảnh đại diện \\ \hline
UC08 & Chọn ngôn ngữ & Khách & Theo hệ thống, tiếng Việt hoặc tiếng Anh \\ \hline
UC09 & Tìm kiếm, lọc khách sạn & Khách & Theo điểm đến, ngày, số khách, giá, điểm \\ \hline
UC10 & Xem chi tiết khách sạn & Khách & Ảnh, tiện nghi, chính sách, đánh giá \\ \hline
UC11 & Xem gợi ý cá nhân hóa & Khách & Hàng ``Đề xuất cho bạn'' trên màn Khám phá \\ \hline
UC12 & Xóa lịch sử gợi ý & Khách & Xóa lượt xem và lượt hiển thị của chính mình \\ \hline
UC13 & Xem thời tiết & Khách, Open-Meteo & Thẻ thời tiết 8 thành phố \\ \hline
UC14 & Lưu khách sạn yêu thích & Khách & Thêm/bỏ khỏi danh sách yêu thích \\ \hline
UC15 & Chọn loại phòng, kiểm tra phòng trống & Khách & Số phòng còn trống cho khoảng ngày \\ \hline
UC16 & Đặt phòng (giữ phòng) & Khách & Tạo đặt phòng chờ thanh toán trong 15 phút \\ \hline
UC17 & Xác nhận trả tại khách sạn & Khách & Chuyển đặt phòng sang đã xác nhận \\ \hline
UC18 & Thanh toán qua payOS & Khách, payOS & Mở trang VietQR; xác nhận bằng webhook \\ \hline
UC19 & Hủy đặt phòng & Khách & Hủy ngay hoặc gửi yêu cầu hủy \\ \hline
UC20 & Xem lịch sử, chi tiết đặt phòng & Khách & Lọc theo trạng thái \\ \hline
UC21 & Đánh giá sau lưu trú & Khách & Một đánh giá cho mỗi khách sạn \\ \hline
UC22 & Xem thông báo & Khách & Đánh dấu đã đọc \\ \hline
UC23 & Gửi yêu cầu hỗ trợ & Khách & Tạo phiếu hỗ trợ \\ \hline
UC24 & Đăng ký làm chủ khách sạn & Khách & Gửi tên doanh nghiệp, mã giấy phép \\ \hline
UC25 & Xem bảng điều khiển, doanh thu & Owner & Thống kê theo khách sạn và theo tháng \\ \hline
UC26 & Quản lý khách sạn & Owner & Tạo, sửa, gửi duyệt, ngừng \\ \hline
UC27 & Quản lý loại phòng, tồn kho & Owner & Giá, sức chứa, số phòng \\ \hline
UC28 & Xem đặt phòng của khách sạn & Owner & Danh sách, tìm kiếm, chi tiết \\ \hline
UC29 & Nhận phòng, trả phòng & Owner & Ghi nhận thời điểm khách đến và đi \\ \hline
UC30 & Xử lý yêu cầu hủy & Owner, Admin & Xác nhận hủy khi chưa có thanh toán trực tuyến \\ \hline
UC31 & Phản hồi đánh giá & Owner & Trả lời đánh giá của khách \\ \hline
UC32 & Duyệt chủ khách sạn & Admin & Chấp nhận hoặc từ chối hồ sơ \\ \hline
UC33 & Duyệt khách sạn & Admin & Đăng hoặc từ chối kèm lý do \\ \hline
UC34 & Quản lý người dùng & Admin & Tìm kiếm, khóa, mở khóa \\ \hline
UC35 & Kiểm duyệt đánh giá & Admin & Ẩn đánh giá, tính lại điểm \\ \hline
UC36 & Giám sát hệ thống & Admin & Thống kê, đặt phòng, thanh toán, nhật ký \\ \hline
UC37 & Trả lời yêu cầu hỗ trợ & Admin & Đóng phiếu hỗ trợ \\ \hline
\end{longtable}
\end{rtable}

\subsection{Biểu đồ ca sử dụng tổng quát}

\Cref{fig:uc-overall} trình bày biểu đồ ca sử dụng tổng quát. Để biểu đồ dễ đọc, các ca sử dụng được gom thành chín nhóm chức năng; mỗi nhóm được tách thành biểu đồ chi tiết ở các mục tiếp theo. Tác nhân chính đứng bên trái, tác nhân phụ là hệ thống bên ngoài đứng bên phải ranh giới hệ thống.

\begin{figure}[H]
\centering
\diagram{uc-overall}
\caption{Biểu đồ ca sử dụng tổng quát của Hotelia}\label{fig:uc-overall}
\end{figure}

\subsection{Biểu đồ ca sử dụng chi tiết}

\paragraph{Nhóm xác thực và tài khoản.} \Cref{fig:uc-auth} chi tiết nhóm xác thực. Đăng nhập Google/Apple \emph{bao gồm} (include) việc xác minh ID token với nhà cung cấp. Đổi mật khẩu và đặt lại mật khẩu \emph{bao gồm} việc tăng \code{tokenVersion} để thu hồi các phiên cũ. Mọi thao tác cần đăng nhập đều mở rộng từ ca làm mới phiên khi access token hết hạn.

\begin{figure}[H]
\centering
\diagram{uc-auth}
\caption{Biểu đồ ca sử dụng nhóm xác thực và tài khoản}\label{fig:uc-auth}
\end{figure}

\paragraph{Nhóm khám phá và gợi ý.} \Cref{fig:uc-explore} cho thấy khách chưa đăng nhập vẫn dùng được tìm kiếm, xem chi tiết và gợi ý khám phá. Gợi ý cá nhân hóa \emph{bao gồm} việc ghi nhận lượt hiển thị và lượt mở chi tiết. Đây là dữ liệu đầu vào cho các lần gợi ý sau.

\begin{figure}[H]
\centering
\diagram{uc-explore}
\caption{Biểu đồ ca sử dụng nhóm khám phá, gợi ý và thời tiết}\label{fig:uc-explore}
\end{figure}

\paragraph{Nhóm đặt phòng và thanh toán.} \Cref{fig:uc-booking} thể hiện đặt phòng \emph{bao gồm} kiểm tra phòng trống. Hai ca thanh toán là \emph{mở rộng} (extend) của đặt phòng, vì khách chọn một trong hai. Thanh toán payOS có cổng thanh toán tham gia với vai trò tác nhân phụ.

\begin{figure}[H]
\centering
\diagram{uc-booking}
\caption{Biểu đồ ca sử dụng nhóm đặt phòng và thanh toán}\label{fig:uc-booking}
\end{figure}

\paragraph{Nhóm chủ khách sạn.} \Cref{fig:uc-owner} gồm quản lý khách sạn, loại phòng, đặt phòng và đánh giá. Mọi ca sử dụng của nhóm này đều \emph{bao gồm} việc kiểm tra quyền sở hữu khách sạn. Kiểm tra này nằm ở tầng dịch vụ, bổ sung cho kiểm tra vai trò.

\begin{figure}[H]
\centering
\diagram{uc-owner}
\caption{Biểu đồ ca sử dụng nhóm chủ khách sạn}\label{fig:uc-owner}
\end{figure}

\paragraph{Nhóm quản trị viên.} \Cref{fig:uc-admin} gồm duyệt, quản lý người dùng, kiểm duyệt và giám sát. Ẩn đánh giá \emph{bao gồm} việc tính lại điểm trung bình của khách sạn.

\begin{figure}[H]
\centering
\diagram{uc-admin}
\caption{Biểu đồ ca sử dụng nhóm quản trị viên}\label{fig:uc-admin}
\end{figure}

\subsection{Đặc tả ca sử dụng}\label{sec:uc-spec}

Các bảng dưới đây đặc tả toàn bộ 37 ca sử dụng ở \cref{tab:uc-list}, theo thứ tự mã. Một số ca sử dụng có quan hệ chặt được đặc tả chung trong một bảng (UC19 với UC30, UC32 với UC33). Các bảng dùng cùng một khuôn dạng gồm tác nhân, tiền điều kiện, luồng chính, luồng thay thế và hậu điều kiện. Biểu đồ tuần tự tương ứng được trình bày ở \cref{sec:sequences}.

\newcommand{\ucspec}[8]{%
\begin{rtable}
\begin{longtable}{|L{3.0cm}|L{11.8cm}|}
\caption{#2}\label{#1}\\
\hline
\thead{Thuộc tính} & \thead{Nội dung} \\ \hline
\endfirsthead
\hline
\thead{Thuộc tính} & \thead{Nội dung} \\ \hline
\endhead
Tác nhân & #3 \\ \hline
Tiền điều kiện & #4 \\ \hline
Luồng chính & #5 \\ \hline
Luồng thay thế & #6 \\ \hline
Hậu điều kiện & #7 \\ \hline
Quy tắc liên quan & #8 \\ \hline
\end{longtable}
\end{rtable}}

\ucspec{tab:uc01}{Đặc tả UC01 -- Đăng ký tài khoản}
{Khách chưa có tài khoản}
{Email chưa được dùng cho tài khoản nào.}
{1. Người dùng nhập email, mật khẩu, họ tên và số điện thoại (không bắt buộc). 2. Ứng dụng gửi \code{POST /auth/register}. 3. \code{ValidationPipe} kiểm tra email hợp lệ, mật khẩu 8--128 ký tự có chữ hoa, chữ thường và chữ số, họ tên 2--100 ký tự, số điện thoại 10--15 chữ số. 4. Máy chủ chuẩn hóa email về chữ thường và kiểm tra trùng. 5. Băm mật khẩu bằng bcrypt, tạo người dùng vai trò USER. 6. Phát hành cặp token như UC02. 7. Ứng dụng lưu token vào Keychain và vào màn Khám phá.}
{3a. Dữ liệu sai định dạng hoặc có trường lạ: 400 kèm thông báo của từng trường. 4a. Email đã đăng ký: 409 ``Email already registered''.}
{Có tài khoản USER mới và người dùng đã đăng nhập.}
{FR-01; NFR-03.}

\ucspec{tab:uc02}{Đặc tả UC02 -- Đăng nhập email -- mật khẩu}
{Khách, chủ khách sạn, quản trị viên}
{Tài khoản tồn tại và đang hoạt động.}
{1. Người dùng nhập email, mật khẩu. 2. Ứng dụng gửi \code{POST /auth/login}. 3. Máy chủ tìm người dùng theo email đã chuẩn hóa chữ thường. 4. So khớp bcrypt. 5. Đặt lại bộ đếm sai, ghi \code{lastLoginAt}. 6. Phát hành access token (15 phút) và refresh token (7 ngày), lưu bản băm SHA-256 của refresh token. 7. Ứng dụng lưu hai token vào Keychain và chuyển sang màn Khám phá.}
{3a. Không tìm thấy email: trả 401 ``Invalid credentials'' (không tiết lộ email có tồn tại hay không). 4a. Sai mật khẩu: tăng bộ đếm; lần thứ 5 khóa tài khoản 15 phút. 4b. Tài khoản đang bị khóa tạm thời: trả 400 ``Account is temporarily locked''. 4c. Tài khoản bị quản trị viên vô hiệu hóa: trả 401 ``Account deactivated''.}
{Người dùng có phiên hợp lệ.}
{NFR-03, NFR-04; \cref{sec:impl-auth}.}

\ucspec{tab:uc03}{Đặc tả UC03 -- Đăng nhập Google/Apple}
{Khách; Google hoặc Apple (IdP)}
{Máy chủ đã cấu hình \code{GOOGLE\_IOS\_CLIENT\_ID} hoặc \code{APPLE\_BUNDLE\_ID}. Với Apple, ứng dụng được ký bằng tài khoản Apple Developer Program.}
{1. Ứng dụng mở luồng đăng nhập của IdP (với Apple, gửi kèm giá trị băm SHA-256 của một nonce ngẫu nhiên). 2. IdP trả ID token. 3. Ứng dụng gửi \code{POST /auth/social} với \code{provider}, \code{idToken} (và \code{nonce} với Apple). 4. Máy chủ xác minh chữ ký RS256 bằng JWKS, \code{iss}, \code{aud}, \code{exp}, email đã xác minh, nonce. 5. Tìm người dùng theo \code{googleSub}/\code{appleSub}; nếu chưa có thì liên kết theo email (chỉ khi email do IdP có thẩm quyền xác nhận) hoặc tạo mới. 6. Phát hành token như UC02.}
{4a. Token sai chữ ký, sai \code{aud} hoặc hết hạn: 401. 4b. Email chưa xác minh: 401. 4c. Nonce Apple không khớp: 401. 1a. Máy chủ chưa cấu hình IdP: 503.}
{Người dùng có phiên hợp lệ; tài khoản được gắn \code{googleSub}/\code{appleSub}.}
{FR-03, FR-04; \cite{oidc,googleidtoken,applesignin}.}

\ucspec{tab:uc04}{Đặc tả UC04 -- Làm mới phiên, đăng xuất}
{Mọi người dùng đã đăng nhập}
{Ứng dụng đang giữ một refresh token.}
{1. Một yêu cầu API trả 401 vì access token (15 phút) hết hạn. 2. \code{AccessTokenRefresher} của ứng dụng gộp các yêu cầu đồng thời thành một lần gọi \code{POST /auth/refresh}. 3. Máy chủ xác minh chữ ký, từ chối token có \code{typ = access}, so refresh token với bản băm đang lưu và kiểm tra tài khoản còn hoạt động. 4. Máy chủ phát hành cặp token mới và thay bản băm cũ (xoay vòng). 5. Ứng dụng lưu cặp mới và gửi lại yêu cầu ban đầu. 6. Khi người dùng chọn đăng xuất, ứng dụng xóa token trong Keychain và dữ liệu người dùng lưu tạm, đồng thời gửi \code{POST /auth/logout} kèm refresh token ở nền (không chờ phản hồi). 7. Máy chủ xác minh refresh token, khớp bản băm thì xóa bản băm, làm phiên không thể làm mới nữa.}
{3a. Refresh token sai, hết hạn, đã bị thay hoặc là access token: 401, ứng dụng kết thúc phiên và về màn đăng nhập. 6a. Mất mạng khi đăng xuất: thiết bị vẫn đăng xuất; refresh token còn hiệu lực tới khi hết hạn. 7a. Token gửi lên là access token, đã hết hạn, giả mạo hoặc đã bị thu hồi: máy chủ vẫn trả 200 và không thay đổi gì, để không lộ thông tin về phiên.}
{Phiên được gia hạn; hoặc token bị xóa khỏi thiết bị và bị thu hồi trên máy chủ.}
{FR-05; \cref{fig:seq-refresh}.}

\ucspec{tab:uc05}{Đặc tả UC05 -- Quên và đặt lại mật khẩu}
{Khách; dịch vụ gửi email}
{Máy chủ đã cấu hình \code{RESET\_EMAIL\_WEBHOOK\_URL} (HTTPS) và \code{RESET\_EMAIL\_WEBHOOK\_TOKEN}.}
{1. Người dùng nhập email và gửi \code{POST /auth/forgot-password}. 2. Máy chủ tìm tài khoản đang hoạt động. 3. Tạo mã ngẫu nhiên 24 byte, lưu bản băm SHA-256 với hạn 15 phút. 4. Gửi mã qua dịch vụ email. 5. Trả cùng một thông báo dù email có tồn tại hay không. 6. Người dùng nhập mã và mật khẩu mới, gửi \code{POST /auth/reset-password}. 7. Trong một giao dịch, máy chủ lấy và xóa mã còn hạn, đổi mật khẩu, tăng \code{tokenVersion} và xóa refresh token.}
{1a. Chưa cấu hình dịch vụ email: 503. 2a. Email không tồn tại hoặc đã yêu cầu trong 60 giây qua: vẫn trả thông báo chung, không gửi email. 4a. Gửi email lỗi: xóa mã vừa tạo, trả 503. 7a. Mã sai hoặc hết hạn: 400.}
{Mật khẩu mới có hiệu lực; mọi phiên cũ bị vô hiệu.}
{FR-06; \cref{fig:seq-password}.}

\ucspec{tab:uc06}{Đặc tả UC06 -- Đổi mật khẩu}
{Mọi người dùng đã đăng nhập}
{Người dùng nhớ mật khẩu hiện tại.}
{1. Người dùng nhập mật khẩu hiện tại và mật khẩu mới (8--72 ký tự). 2. Ứng dụng gửi \code{POST /account/password}. 3. Máy chủ so khớp mật khẩu hiện tại bằng bcrypt. 4. Cập nhật có điều kiện theo bản băm cũ: đặt bản băm mới, tăng \code{tokenVersion}, xóa refresh token. 5. Ứng dụng đăng xuất và yêu cầu đăng nhập lại.}
{3a. Sai mật khẩu hiện tại: 400. 3b. Mật khẩu mới trùng mật khẩu cũ: 400. 4a. Mật khẩu vừa bị đổi bởi yêu cầu khác: 400, đề nghị thử lại.}
{Access token cũ bị \code{JwtStrategy} từ chối do lệch \code{tokenVersion}.}
{FR-07.}

\ucspec{tab:uc07}{Đặc tả UC07 -- Xem, cập nhật hồ sơ}
{Mọi người dùng đã đăng nhập}
{Có phiên hợp lệ.}
{1. Màn Tài khoản gọi \code{GET /users/me}. 2. Người dùng sửa họ tên, số điện thoại hoặc ảnh đại diện. 3. Ứng dụng gửi \code{PUT /users/me}. 4. \code{ValidationPipe} chỉ nhận ba trường này. 5. Máy chủ lưu và trả hồ sơ không có \code{passwordHash}.}
{4a. Số điện thoại sai định dạng hoặc họ tên quá 100 ký tự: 400. 4b. Gửi thêm trường như \code{role} hay \code{email}: 400 (ca A11).}
{Hồ sơ được cập nhật; email và vai trò không đổi.}
{FR-08; \cref{fig:seq-profile}.}

\ucspec{tab:uc08}{Đặc tả UC08 -- Chọn ngôn ngữ}
{Người dùng ứng dụng iOS}
{Không có.}
{1. Người dùng mở phần Ngôn ngữ trong màn Tài khoản. 2. Chọn ``Theo hệ thống'', ``Tiếng Việt'' hoặc ``English''. 3. Giá trị được lưu vào \code{UserDefaults} qua \code{@AppStorage("appLanguage")}. 4. \code{DuLichApp} đổi \code{locale} và dựng lại cây giao diện bằng \code{.id(appLanguage)}.}
{2a. Chọn ``Theo hệ thống'': ứng dụng theo ngôn ngữ đang đặt trên iOS.}
{Mọi màn hình hiển thị theo ngôn ngữ mới; lựa chọn được giữ sau khi mở lại ứng dụng. Nội dung dữ liệu (mô tả khách sạn) vẫn là tiếng Việt.}
{NFR-13; \cref{fig:seq-profile}.}

\ucspec{tab:uc09}{Đặc tả UC09 -- Tìm kiếm, lọc khách sạn}
{Khách (không bắt buộc đăng nhập)}
{Không có.}
{1. Khách nhập điểm đến và/hoặc chọn ngày, số khách, bộ lọc. 2. Ứng dụng gửi \code{GET /hotels} với tham số tương ứng. 3. Máy chủ chỉ xét khách sạn \code{PUBLISHED}. 4. Nếu có bộ lọc phòng, máy chủ tìm loại phòng đang bán thỏa sức chứa và khoảng giá, rồi với khoảng ngày đã chọn, chỉ giữ loại phòng còn ít nhất một phòng trống ở \emph{mọi} đêm. 5. Lọc theo điểm đến (thành phố, quận, tên), điểm tối thiểu và tiện nghi. 6. Sắp xếp, phân trang (tối đa 100), trả kết quả.}
{1a. Chỉ có ngày nhận hoặc chỉ có ngày trả: 400. 1b. Giá tối thiểu lớn hơn giá tối đa: 400. 1c. Ngày nhận ở quá khứ hoặc kỳ lưu trú ngoài 1--90 đêm: 400.}
{Danh sách khách sạn mà khách thực sự đặt được với điều kiện đã chọn.}
{BR-01, BR-02, BR-03; \cref{sec:search-design}.}

\ucspec{tab:uc10}{Đặc tả UC10 -- Xem chi tiết khách sạn}
{Khách (không bắt buộc đăng nhập)}
{Khách sạn ở trạng thái PUBLISHED.}
{1. Người dùng chọn một khách sạn. 2. Ứng dụng gọi \code{GET /hotels/:id} và \code{GET /reviews/hotel/:id}. 3. Máy chủ trả ảnh, mô tả, tiện nghi, giờ nhận/trả phòng, điểm đánh giá đã xác minh; điểm biên tập (nếu có) được trả ở trường riêng \code{editorialRating}. 4. Ứng dụng ghi một lượt mở chi tiết qua \code{POST /recommendations/views} để phục vụ gợi ý. 5. Người dùng xem loại phòng (UC15) hoặc lưu yêu thích (UC14).}
{2a. Khách sạn không tồn tại hoặc chưa PUBLISHED: 404. 3a. Khách sạn thuộc danh mục mẫu: \code{bookingEnabled = false}, chỉ xem được, không đặt được (BR-10). 4a. Ghi lượt xem lỗi: bỏ qua, không ảnh hưởng màn hình.}
{Lượt xem được lưu (mỗi khách sạn một bản ghi cho mỗi tác nhân).}
{FR-10.}

\ucspec{tab:uc11}{Đặc tả UC11 -- Xem gợi ý cá nhân hóa}
{Khách (đăng nhập hoặc chưa)}
{Ứng dụng có \code{sessionId} UUID cho thiết bị.}
{1. Màn Khám phá gửi \code{GET /recommendations} với \code{sessionId}, điểm đến và bộ lọc hiện tại; gửi kèm access token nếu có. 2. Máy chủ xác định tác nhân: \code{user:<id>} nếu đăng nhập, ngược lại \code{session:<uuid>}. 3. Đọc lượt xem 90 ngày của tác nhân và đặt phòng đã xác nhận của người dùng. 4. Lọc ứng viên theo điểm đến và bộ lọc phòng. 5. Tính 16 đặc trưng, chấm điểm bằng LightGBM. 6. Trả danh sách, kèm \code{ranking.mode}, \code{version}, \code{modelUsed}. 7. Ứng dụng gửi \code{POST /recommendations/impressions} với thứ tự hiển thị; khi khách mở chi tiết, gửi \code{POST /recommendations/views}.}
{5a. Mô hình không nạp được hoặc lỗi khi chấm điểm: dùng bộ xếp hạng theo nội dung, \code{modelUsed = false}. 3a. Chưa có lịch sử: chế độ khám phá. 7a. Gửi lại lượt hiển thị cùng \code{impressionId}: không ghi đè.}
{Khách thấy danh sách khách sạn đặt được, xếp theo mức phù hợp ước lượng.}
{BR-21 đến BR-24; \cref{sec:rec-design}.}

\ucspec{tab:uc12}{Đặc tả UC12 -- Xóa lịch sử gợi ý}
{Khách}
{Có \code{sessionId} của thiết bị; có thể đã đăng nhập.}
{1. Trong màn Tài khoản, người dùng chọn ``Xóa lịch sử gợi ý'' và xác nhận. 2. Ứng dụng gửi \code{DELETE /recommendations/views?sessionId=\ldots}. 3. Máy chủ xác định tác nhân: \code{user:<id>} nếu có token, ngược lại \code{session:<id>}. 4. Xóa mọi lượt xem và lượt hiển thị của tác nhân đó. 5. Trả \code{\{cleared: true\}}; ứng dụng báo đã xóa.}
{1a. Người dùng hủy ở hộp xác nhận: không có thay đổi. 2a. Lỗi mạng: ứng dụng báo không xóa được và đề nghị thử lại.}
{Gợi ý quay về chế độ khám phá cho tới khi có lịch sử mới.}
{FR-13; NFR-14.}

\ucspec{tab:uc13}{Đặc tả UC13 -- Xem thời tiết}
{Khách; Open-Meteo}
{Không có.}
{1. Màn Khám phá yêu cầu dự báo. 2. \code{WeatherService} đọc bộ đệm trong \code{UserDefaults}. 3. Nếu bộ đệm dưới 30 phút: hiển thị ngay. 4. Nếu không: gọi Open-Meteo một lần cho 8 thành phố (thời tiết hiện tại và 4 ngày tới, múi giờ Asia/Ho\_Chi\_Minh). 5. Lưu bộ đệm và hiển thị thẻ thời tiết. 6. Người dùng mở một thành phố để xem dự báo chi tiết.}
{4a. Lỗi mạng hoặc quá 12 giây: dùng bộ đệm dưới 6 giờ nếu có, nếu không thì hiển thị trạng thái lỗi có nút thử lại.}
{Thẻ thời tiết hiển thị; không có dữ liệu nào gửi về máy chủ Hotelia.}
{FR-14; \cref{fig:seq-weather}.}

\ucspec{tab:uc14}{Đặc tả UC14 -- Lưu khách sạn yêu thích}
{Khách đã đăng nhập}
{Khách sạn ở trạng thái PUBLISHED.}
{1. Người dùng bấm biểu tượng trái tim. 2. Nếu chưa lưu, ứng dụng gửi \code{PUT /account/favorites/:id}; máy chủ \code{upsert} bản ghi có khóa \code{userId:hotelId}. 3. Nếu đã lưu, ứng dụng gửi \code{DELETE /account/favorites/:id}. 4. Ứng dụng cập nhật danh sách sau khi máy chủ trả thành công. 5. Màn Yêu thích tải danh sách qua \code{GET /account/favorites}.}
{1a. Bấm lặp khi yêu cầu đang chạy: bị bỏ qua. 2a. Khách sạn không tồn tại hoặc chưa PUBLISHED: 404. 4a. Yêu cầu lỗi: biểu tượng giữ nguyên, hiển thị thông báo lỗi.}
{Danh sách yêu thích phản ánh đúng trạng thái trên máy chủ.}
{FR-15; \cref{fig:seq-favorites}.}

\ucspec{tab:uc15}{Đặc tả UC15 -- Chọn loại phòng, kiểm tra phòng trống}
{Khách (không bắt buộc đăng nhập)}
{Khách sạn PUBLISHED có loại phòng đang bán.}
{1. Người dùng chọn ngày nhận, ngày trả và số khách. 2. Ứng dụng gọi \code{GET /room-types/hotel/:hotelId?checkIn\&checkOut}. 3. Máy chủ kiểm tra khoảng ngày (BR-01, BR-02). 4. Đọc các đặt phòng chồng lấn đang chiếm tồn kho (BR-04). 5. Tính số phòng trống nhỏ nhất qua các đêm (BR-03). 6. Ứng dụng hiển thị giá, sức chứa và số phòng còn trống của từng loại phòng.}
{2a. Không có ngày: trả loại phòng không kèm số phòng trống. 3a. Chỉ có một trong hai ngày, ngày sai định dạng, ở quá khứ hoặc quá 90 đêm: 400. 6a. Loại phòng hết phòng hoặc không đủ sức chứa: ứng dụng không cho chọn.}
{Người dùng chọn được loại phòng để đặt (UC16).}
{FR-16; BR-01 -- BR-04; \cref{fig:seq-availability}.}

\ucspec{tab:uc16}{Đặc tả UC16 -- Đặt phòng (giữ phòng)}
{Khách đã đăng nhập}
{Khách sạn \code{PUBLISHED}, không thuộc danh mục chỉ xem; loại phòng đang bán.}
{1. Khách chọn loại phòng, ngày, số khách, nhập thông tin liên hệ. 2. Ứng dụng sinh UUID làm \code{Idempotency-Key} cho nội dung này và gửi \code{POST /bookings}. 3. Máy chủ tính băm SHA-256 của nội dung. 4. Nếu đã có đặt phòng với cùng khóa và cùng băm: trả lại đặt phòng đó. 5. Mở giao dịch; tăng \code{inventoryVersion} của loại phòng. 6. Kiểm tra sức chứa, giá. 7. Đếm phòng trống theo từng đêm trong số đặt phòng đang chiếm tồn kho. 8. Tạo đặt phòng \code{PENDING\_PAYMENT} với \code{holdExpiresAt} = hiện tại + 15 phút; ghi thông báo. 9. Xác nhận giao dịch, trả 201.}
{2a. Thiếu khóa hoặc khóa không phải UUID: 400. 4a. Cùng khóa nhưng nội dung khác: 409. 6a. Vượt sức chứa: 400. 7a. Hết phòng: 400 ``Loại phòng đã hết trong thời gian đã chọn''. 5a. Xung đột ghi do yêu cầu đồng thời: trình điều khiển chạy lại giao dịch; lần chạy lại thấy tồn kho mới. E. Mất kết nối trước khi nhận phản hồi: ứng dụng giữ nguyên khóa; khách bấm ``Kiểm tra lại'' để gửi lại cùng khóa.}
{Một đặt phòng duy nhất giữ tồn kho trong 15 phút.}
{BR-04 đến BR-10; \cref{sec:booking-rules}.}

\ucspec{tab:uc17}{Đặc tả UC17 -- Xác nhận trả tại khách sạn}
{Khách đã đăng nhập}
{Đặt phòng của chính người dùng ở trạng thái PENDING\_PAYMENT và còn hạn giữ.}
{1. Ở màn chi tiết đặt phòng, người dùng chọn trả tại khách sạn. 2. Ứng dụng gửi \code{POST /bookings/:id/pay-at-hotel}. 3. Máy chủ kiểm tra người gọi là người đặt. 4. Cập nhật có điều kiện sang CONFIRMED, \code{paymentMethod = PAY\_AT\_HOTEL}. 5. Ứng dụng hiển thị trạng thái đã xác nhận.}
{3a. Không phải người đặt: 403. 4a. Đã CONFIRMED: trả lại đặt phòng hiện tại. 4b. Không còn PENDING\_PAYMENT hoặc hạn giữ vừa hết: 400.}
{Đặt phòng CONFIRMED, vẫn chiếm tồn kho; không có bản ghi thanh toán.}
{FR-19; BR-04, BR-05; \cref{fig:seq-pay-at-hotel}.}

\ucspec{tab:uc18}{Đặc tả UC18 -- Thanh toán qua payOS}
{Khách; payOS}
{Đặt phòng \code{PENDING\_PAYMENT} của chính khách, chưa hết hạn giữ; máy chủ đã cấu hình khóa payOS.}
{1. Khách chọn ``Thanh toán qua payOS''. 2. Ứng dụng gửi \code{POST /payments} với \code{bookingId}. 3. Máy chủ tạo bản ghi thanh toán \code{PROCESSING} với \code{orderCode} duy nhất, số tiền lấy từ đặt phòng (không lấy từ máy khách), hạn bằng hạn giữ phòng. 4. Gọi payOS tạo liên kết; kiểm tra đặt phòng vẫn đang giữ; chuyển thanh toán sang \code{PENDING}, trả \code{checkoutUrl}. 5. Ứng dụng mở trang thanh toán. 6. payOS gửi webhook; máy chủ xác minh chữ ký, đối chiếu số tiền, trong một giao dịch chuyển thanh toán sang \code{COMPLETED} và đặt phòng sang \code{CONFIRMED}, ghi thông báo.}
{2a. Đã có liên kết đang chờ: trả lại liên kết cũ (không tạo đơn thứ hai). 3a. payOS chưa cấu hình: 503. 4a. payOS không phản hồi trong 10 giây: thanh toán \code{FAILED}, trả 502, khách thử lại được. 4b. Giữ phòng đã hết khi liên kết vừa tạo: hủy liên kết, 400. 6a. Chữ ký sai: 400, không đổi dữ liệu. 6b. Số tiền không khớp: hủy giao dịch cơ sở dữ liệu, không đổi trạng thái. 6c. Tiền về sau khi hết giữ phòng: đặt phòng chuyển \code{CANCEL\_REQUESTED} với lý do \code{PAYMENT\_AFTER\_HOLD\_EXPIRED} để xử lý hoàn tiền thủ công. 6d. Webhook gửi lại: bỏ qua vì đã \code{COMPLETED}.}
{Đặt phòng được xác nhận chỉ khi tiền đã về và khớp số tiền.}
{BR-11 đến BR-14; \cref{sec:payment-design}.}

\ucspec{tab:uc19}{Đặc tả UC19 -- Hủy đặt phòng và UC30 -- Xử lý yêu cầu hủy}
{Khách; chủ khách sạn hoặc quản trị viên}
{Đặt phòng thuộc khách (UC19) hoặc thuộc khách sạn của chủ khách sạn (UC30).}
{UC19: 1. Khách chọn hủy, nhập lý do. 2. Nếu đặt phòng đang \code{PENDING\_PAYMENT}: chuyển \code{CANCELLED}, giải phóng tồn kho, hủy liên kết payOS đang chờ (nếu có). 3. Nếu đặt phòng \code{CONFIRMED}: chuyển \code{CANCEL\_REQUESTED}; tồn kho \emph{vẫn bị giữ}. UC30: 4. Chủ khách sạn hoặc quản trị viên xác nhận hủy; nếu không có thanh toán trực tuyến đã hoàn tất, đặt phòng chuyển \code{CANCELLED}.}
{2a. Đặt phòng đã ở trạng thái đích: trả lại nguyên trạng (lũy đẳng). 2b. Đặt phòng đã nhận phòng hoặc kết thúc: 400. 4a. Có thanh toán payOS đã hoàn tất: 400 ``Cần xử lý hoàn tiền với nhà cung cấp trước'', vì hoàn tiền tự động chưa được triển khai.}
{Tồn kho chỉ được giải phóng khi đặt phòng thật sự bị hủy.}
{BR-15, BR-16.}

\ucspec{tab:uc20}{Đặc tả UC20 -- Xem lịch sử, chi tiết đặt phòng}
{Khách đã đăng nhập}
{Có phiên hợp lệ.}
{1. Người dùng mở tab Đặt phòng. 2. Ứng dụng gọi \code{GET /bookings}. 3. Máy chủ chuyển các giữ phòng quá hạn của người dùng sang EXPIRED, rồi trả danh sách mới nhất trước. 4. Ứng dụng chia thẻ Tất cả, Sắp tới, Đã xong, Đã hủy. 5. Người dùng chọn một đặt phòng; ứng dụng gọi \code{GET /bookings/:id} để xem chi tiết và các thao tác còn được phép (thanh toán, hủy, đánh giá).}
{5a. Đặt phòng không thuộc người dùng: 403. 5b. Mã không tồn tại: 404. 4a. Danh sách rỗng: hiển thị trạng thái trống.}
{Người dùng thấy trạng thái hiện tại của mọi đặt phòng.}
{FR-23; BR-05; \cref{fig:seq-history}.}

\ucspec{tab:uc21}{Đặc tả UC21 -- Đánh giá sau lưu trú}
{Khách}
{Khách có ít nhất một đặt phòng \code{CHECKED\_OUT} hoặc \code{COMPLETED} tại khách sạn; chưa đánh giá khách sạn này.}
{1. Khách chọn số sao (1--5), tiêu đề, nội dung. 2. Trong một giao dịch: kiểm tra điều kiện, tạo đánh giá, tính lại \code{averageRating} và \code{reviewCount} của khách sạn từ các đánh giá đang hiển thị. 3. Chủ khách sạn có thể phản hồi (UC31); phản hồi hiển thị dưới đánh giá trên ứng dụng.}
{1a. Chưa lưu trú: 403 ``Chỉ có thể đánh giá sau khi hoàn tất lưu trú''. 1b. Đã đánh giá: 403.}
{Điểm đánh giá đã xác minh của khách sạn được cập nhật.}
{BR-19, BR-20.}

\ucspec{tab:uc22}{Đặc tả UC22 -- Xem thông báo}
{Khách đã đăng nhập}
{Có thông báo do hệ thống tạo khi giữ phòng hoặc khi thanh toán payOS thành công.}
{1. Người dùng mở màn Thông báo. 2. Ứng dụng gọi \code{GET /notifications?page}, nhận tối đa 20 thông báo mỗi trang và số chưa đọc. 3. Người dùng mở một thông báo; ứng dụng gửi \code{POST /notifications/:id/read}. 4. Người dùng có thể chọn đọc tất cả qua \code{POST /notifications/read-all}.}
{3a. Mã thông báo không thuộc người dùng: không có thay đổi.}
{Trạng thái đã đọc được lưu trên máy chủ.}
{FR-25.}

\ucspec{tab:uc23}{Đặc tả UC23 -- Gửi yêu cầu hỗ trợ}
{Khách đã đăng nhập}
{Có phiên hợp lệ.}
{1. Người dùng mở màn Hỗ trợ; ứng dụng gọi \code{GET /account/support} để hiện 100 phiếu gần nhất. 2. Người dùng nhập tiêu đề (3--160 ký tự) và nội dung (10--4.000 ký tự). 3. Ứng dụng gửi \code{POST /account/support}. 4. Máy chủ tạo phiếu trạng thái OPEN.}
{3a. Tiêu đề hoặc nội dung quá ngắn: 400.}
{Phiếu chờ quản trị viên trả lời (UC37).}
{FR-37.}

\ucspec{tab:uc24}{Đặc tả UC24 -- Đăng ký làm chủ khách sạn}
{Khách đã đăng nhập}
{Tài khoản có vai trò USER và chưa có hồ sơ đang chờ duyệt.}
{1. Người dùng nhập tên doanh nghiệp và mã giấy phép kinh doanh (không bắt buộc). 2. Ứng dụng gửi \code{POST /owners/apply}. 3. Máy chủ lưu thông tin doanh nghiệp và đặt \code{ownerStatus = PENDING}. 4. Hồ sơ xuất hiện trong danh sách chờ duyệt của quản trị viên.}
{3a. Tài khoản đã là OWNER/ADMIN hoặc đang có hồ sơ chờ duyệt: 400.}
{Vai trò vẫn là USER cho tới khi được duyệt (UC32).}
{FR-26; \cref{fig:seq-approve}.}

\ucspec{tab:uc25}{Đặc tả UC25 -- Xem bảng điều khiển, doanh thu}
{Chủ khách sạn (hoặc quản trị viên)}
{Đăng nhập với vai trò OWNER hoặc ADMIN.}
{1. Chủ khách sạn mở trang Tổng quan. 2. Bảng điều khiển gọi \code{GET /owners/dashboard}. 3. Máy chủ lấy khách sạn thật của chủ, đếm khách sạn PUBLISHED và chờ duyệt, tổng số đặt phòng. 4. Tính tổng doanh thu, doanh thu hôm nay và tháng này theo giờ Việt Nam từ các đặt phòng đã xác nhận trở đi, cùng tỷ lệ phòng đang có khách. 5. Trang Doanh thu gọi \code{GET /owners/revenue?year} để lấy doanh thu theo tháng.}
{2a. Vai trò USER: 403. 3a. Chủ chưa có khách sạn: các chỉ số bằng 0.}
{Chủ khách sạn thấy số liệu của riêng mình.}
{FR-30; \cref{fig:seq-owner-dashboard}.}

\ucspec{tab:uc26}{Đặc tả UC26 -- Quản lý khách sạn}
{Chủ khách sạn}
{Vai trò OWNER.}
{1. Chủ khách sạn nhập tên, mô tả, địa chỉ, thành phố, hạng sao, tiện nghi, giờ nhận/trả phòng, ảnh. 2. \code{POST /hotels} tạo khách sạn ở trạng thái DRAFT. 3. Chủ sửa thông tin qua \code{PUT /hotels/:id}. 4. Chủ gửi duyệt qua \code{POST /hotels/:id/submit}, trạng thái chuyển PENDING\_APPROVAL. 5. Sau khi được duyệt (UC33), khách sạn xuất hiện trong tìm kiếm. 6. Chủ có thể ngừng khách sạn qua \code{DELETE /hotels/:id} (chuyển SUSPENDED, không xóa dữ liệu).}
{2a, 3a, 4a. Khách sạn không thuộc chủ hoặc thuộc danh mục mẫu: 403. 4b. Gửi duyệt khi không ở DRAFT hoặc REJECTED: 403. 4c. Bị từ chối: trạng thái REJECTED kèm lý do, chủ sửa rồi gửi lại.}
{Khách sạn đi theo máy trạng thái ở \cref{fig:state-hotel-payment}.}
{FR-27; \cref{fig:seq-approve}.}

\ucspec{tab:uc27}{Đặc tả UC27 -- Quản lý loại phòng, tồn kho}
{Chủ khách sạn}
{Khách sạn thuộc chủ và không thuộc danh mục mẫu.}
{1. Chủ tạo loại phòng qua \code{POST /room-types/hotel/:hotelId} với tên, giá, sức chứa, số phòng. 2. Chủ sửa qua \code{PUT /room-types/:id}. 3. Trong một giao dịch, máy chủ tăng \code{inventoryVersion}, kiểm tra quyền sở hữu. 4. Nếu giảm \code{totalRooms}, đếm số phòng đã bán tại mỗi mốc ngày tương lai. 5. Lưu và commit. 6. Ngừng bán qua \code{DELETE /room-types/:id} (đặt \code{isActive = false} bằng cùng luồng cập nhật).}
{1a, 3a. Khách sạn không thuộc chủ: 403; thuộc danh mục mẫu: 403. 4a. Số phòng mới nhỏ hơn số đã bán ở một ngày nào đó: hủy giao dịch, 400.}
{Tồn kho mới không làm bán trùng các đặt phòng đã có.}
{FR-28; BR-08, BR-25; \cref{fig:seq-inventory}.}

\ucspec{tab:uc28}{Đặc tả UC28 -- Xem đặt phòng của khách sạn}
{Chủ khách sạn}
{Vai trò OWNER.}
{1. Chủ mở trang Đặt phòng. 2. Bảng điều khiển gọi \code{GET /owners/bookings} (mọi khách sạn) hoặc \code{GET /owners/hotels/:hotelId/bookings} (một khách sạn, kèm họ tên, email, số điện thoại của khách). 3. Máy chủ chuyển các giữ phòng quá hạn sang EXPIRED rồi trả danh sách. 4. Chủ lọc, tìm kiếm và mở chi tiết qua \code{GET /bookings/:id}.}
{2a. Khách sạn không thuộc chủ: 404 (ca A7). 4a. Chi tiết đặt phòng ở khách sạn khác: 403 (ca A13).}
{Chủ chỉ thấy đặt phòng của khách sạn mình.}
{FR-29.}

\ucspec{tab:uc29}{Đặc tả UC29 -- Nhận phòng, trả phòng}
{Chủ khách sạn}
{Đặt phòng thuộc khách sạn của chủ khách sạn.}
{1. Chủ khách sạn mở danh sách đặt phòng, chọn ``Nhận phòng''. 2. Máy chủ kiểm tra quyền sở hữu, kiểm tra chuyển trạng thái \code{CONFIRMED} $\rightarrow$ \code{CHECKED\_IN} hợp lệ, kiểm tra ngày hiện tại (giờ Việt Nam) nằm trong $[d_{\text{in}}, d_{\text{out}})$. 3. Cập nhật có điều kiện theo trạng thái cũ, ghi \code{checkedInAt}. 4. Khi khách đi, chọn ``Trả phòng'': \code{CHECKED\_IN} $\rightarrow$ \code{CHECKED\_OUT}, ghi \code{checkedOutAt}.}
{2a. Không phải chủ khách sạn: 403. 2b. Chưa đến ngày nhận phòng hoặc đã qua ngày trả phòng: 400. 2c. Trạng thái không cho phép: 400. 3a. Bấm lại khi đã nhận phòng: trả lại nguyên trạng.}
{Khách được phép đánh giá sau khi đặt phòng ở trạng thái \code{CHECKED\_OUT}.}
{BR-17, BR-18.}

\ucspec{tab:uc31}{Đặc tả UC31 -- Phản hồi đánh giá}
{Chủ khách sạn}
{Đánh giá thuộc một khách sạn của chủ.}
{1. Chủ mở trang Đánh giá; bảng điều khiển gọi \code{GET /reviews/owner}. 2. Chủ nhập phản hồi và gửi \code{POST /reviews/:id/reply}. 3. Máy chủ kiểm tra quyền sở hữu khách sạn, lưu \code{ownerReply}. 4. Ứng dụng iOS hiển thị khối ``Phản hồi từ khách sạn'' dưới đánh giá.}
{3a. Đánh giá không tồn tại: 404. 3b. Khách sạn không thuộc chủ: 403. 2a. Gửi lại: phản hồi mới thay phản hồi cũ.}
{Phản hồi hiển thị công khai cùng đánh giá.}
{FR-31.}

\ucspec{tab:uc33}{Đặc tả UC32, UC33 -- Duyệt chủ khách sạn và khách sạn}
{Quản trị viên}
{Có hồ sơ chủ khách sạn \code{PENDING} hoặc khách sạn \code{PENDING\_APPROVAL}.}
{UC32: 1. Quản trị viên mở danh sách chủ khách sạn chờ duyệt. 2. Chọn ``Duyệt'': người dùng chuyển vai trò sang OWNER, \code{ownerStatus = APPROVED}. UC33: 3. Mở danh sách khách sạn chờ duyệt. 4. Chọn ``Duyệt'': khách sạn chuyển \code{PUBLISHED}, xuất hiện trong tìm kiếm.}
{2a. Từ chối: \code{ownerStatus = REJECTED}. 4a. Từ chối khách sạn: bắt buộc nhập lý do; khách sạn chuyển \code{REJECTED}, chủ khách sạn sửa và gửi duyệt lại. 4b. Khách sạn không ở trạng thái chờ duyệt: 400.}
{Chỉ khách sạn đã duyệt mới được tìm thấy và đặt.}
{BR-25.}

\ucspec{tab:uc34}{Đặc tả UC34 -- Quản lý người dùng}
{Quản trị viên}
{Vai trò ADMIN.}
{1. Quản trị viên mở trang Người dùng; bảng điều khiển gọi \code{GET /admin/users?page\&search\&role}. 2. Máy chủ tìm theo họ tên hoặc email (đã thoát ký tự đặc biệt), lọc vai trò, trả tối đa 100 bản ghi mỗi trang, không kèm mật khẩu và refresh token. 3. Quản trị viên khóa qua \code{POST /admin/users/:id/suspend} hoặc mở khóa qua \code{POST /admin/users/:id/activate}. 4. Thao tác được ghi vào nhật ký.}
{3a. Người dùng không tồn tại: 404. 3b. Vai trò khác ADMIN: 403 (ca A3).}
{Tài khoản bị khóa không đăng nhập được và mọi yêu cầu kế tiếp bị 401.}
{FR-34; \cref{fig:seq-admin-moderate}.}

\ucspec{tab:uc35}{Đặc tả UC35 -- Kiểm duyệt đánh giá}
{Quản trị viên}
{Vai trò ADMIN.}
{1. Quản trị viên xem 200 đánh giá mới nhất qua \code{GET /admin/reviews}. 2. Chọn ẩn một đánh giá vi phạm, gửi \code{POST /reviews/:id/hide}. 3. Trong một giao dịch, máy chủ tăng \code{reviewVersion} của khách sạn, đặt \code{isVisible = false} và tính lại số đánh giá, điểm trung bình từ các đánh giá còn hiển thị.}
{3a. Đánh giá không tồn tại: 404.}
{Đánh giá không còn hiển thị; điểm của khách sạn được cập nhật.}
{FR-35; BR-20; \cref{fig:seq-admin-moderate}.}

\ucspec{tab:uc36}{Đặc tả UC36 -- Giám sát hệ thống}
{Quản trị viên}
{Vai trò ADMIN.}
{1. Trang Tổng quan gọi \code{GET /admin/stats}: số người dùng, chủ khách sạn, khách sạn PUBLISHED, đặt phòng, đánh giá đang hiển thị. 2. Các trang Khách sạn, Đặt phòng, Thanh toán, Đánh giá gọi \code{GET /admin/hotels}, \code{/admin/bookings}, \code{/admin/payments}, \code{/admin/reviews}. 3. Trang Nhật ký gọi \code{GET /admin/audit-logs}, trả 200 bản ghi mới nhất gồm người thực hiện, vai trò, phương thức, đường dẫn và thời điểm.}
{1a. Vai trò khác ADMIN: 403 (ca A2).}
{Quản trị viên nắm được tình trạng hệ thống; nhật ký không chứa nội dung yêu cầu.}
{FR-36; \cref{fig:seq-admin-stats}.}

\ucspec{tab:uc37}{Đặc tả UC37 -- Trả lời yêu cầu hỗ trợ}
{Quản trị viên}
{Có phiếu hỗ trợ do khách gửi (UC23).}
{1. Quản trị viên xem 200 phiếu mới nhất qua \code{GET /account/support/all}. 2. Nhập nội dung trả lời (2--4.000 ký tự) và gửi \code{POST /account/support/:id/reply}. 3. Máy chủ lưu câu trả lời, đổi trạng thái sang RESOLVED và ghi thời điểm trả lời. 4. Khách thấy câu trả lời trong màn Hỗ trợ.}
{3a. Phiếu không tồn tại: 404. 2a. Vai trò khác ADMIN: 403.}
{Phiếu được đóng.}
{FR-37.}

